% ═══════════════════════════════════════════════════════════════════════
% ឈ្មោះឯកសារ: slide_wrapper.tex
%
% ការពិពណ៌នា: ទម្រង់ Standalone សម្រាប់ចងក្រងសំណួរ/លំហាត់នីមួយៗ \begin{ex}
%             ចេញជារូបភាពស្លាយ — រាល់លំហាត់នីមួយៗត្រូវបានរុំក្នុងស៊ុម tcolorbox + tikz។
%
% Placeholder ខ្លឹមសារ: SLIDE_BODY (ជំនួសដោយស្វ័យប្រវត្តិក្នង \begin{document})
% ═══════════════════════════════════════════════════════════════════════

\documentclass[tikz]{standalone}

\usepackage[hidelinks]{hyperref}
\usepackage{amsmath,amssymb,yhmath,mathrsfs}
\usepackage{tikz}
\usepackage{tkz-euclide}
\usepackage{tikz-3dplot,tkz-tab}
\usetikzlibrary{shapes.geometric,arrows,snakes,shadows,backgrounds,decorations.text,shapes}
\usetikzlibrary{calc,intersections,angles,patterns,backgrounds,positioning}
\usetikzlibrary{decorations.pathmorphing,fadings,shadings}
\usepackage{bbding,pifont}
\usepackage{enumerate}
\usepackage{makecell}
\usepackage{framed}
\usepackage[many]{tcolorbox}
\tcbuselibrary{skins,breakable,theorems}
\colorlet{tcbcol@back}{tcbcolback}
\colorlet{tcbcol@frame}{tcbcolframe}
\usepackage{esvect}
\def\vec{\vv}
\def\overrightarrow{\vv}

% កញ្ចប់ ex_test និងការកំណត់ពុម្ពអក្សរខ្មែរ
\usepackage[dethi]{ex_test}
\usepackage{iftex}
\ifXeTeX
  \usepackage{fontspec}
  \IfFontExistsTF{Khmer OS}{\setmainfont{Khmer OS}}{%
    \IfFontExistsTF{Noto Sans Khmer}{\setmainfont{Noto Sans Khmer}}{%
      \IfFontExistsTF{Khmer Sangam MN}{\setmainfont{Khmer Sangam MN}}{}%
    }%
  }
\fi
\renewcommand{\nameex}{\color{violet}លំហាត់}
\renewcommand{\namebt}{\color{blue}លំហាត់}
\renewcommand{\namevd}{\color{red}ឧទាហរណ៍}

% ── ពណ៌ (Color Palette) ──────────────────────────────────────────────
\definecolor{FrameNavy}   {RGB}{20,  60, 130}
\definecolor{BgSlide}     {RGB}{240, 245, 255}
\definecolor{AccentGold}  {RGB}{200, 150,  20}

% ── កំណត់ទំហំរូបភាព tikzpicture ក្នុងខ្លឹមសារលំហាត់ ──────────────────
% រូប TikZ នីមួយៗត្រូវបានចាប់យកដាក់ក្នុង lrbox ហើយ \resizebox នឹងបង្រួម/ពង្រីកមកទំហំដែលបានកំណត់។
\newlength{\SlideTikzPictureWidth}
\newlength{\SlideTikzPictureHeight}
\setlength{\SlideTikzPictureWidth}{5cm}
\setlength{\SlideTikzPictureHeight}{3cm}
\newsavebox{\SlideTikzBox}

% ── កំណត់ឡើងវិញនូវ ex: លំហាត់នីមួយៗ = tikzpicture + tcolorbox ─────────
\RenewEnviron{ex}{%
  \begin{tikzpicture}
    \node[scale=1] {%
      \begin{tcolorbox}[
        width=16cm,
        enhanced,
        %colback=BgSlide,
        colframe=FrameNavy,
        arc=7pt,
        boxrule=1.5pt,
        left=12pt, right=12pt, top=10pt, bottom=10pt,
        drop shadow={black!25, shadow xshift=2pt, shadow yshift=-2pt},
        fonttitle=\bfseries\color{FrameNavy},
      ]
      % Group {} កំណត់ដែនកំណត់នៃការ redefine៖ គ្រប់ tikzpicture ខាងក្នុង \BODY នឹង
      % ត្រូវបានរុំដោយស្វ័យប្រវត្តិតាមរយៈ savebox + resizebox
      {%
        \let\OrigTikzPicture\tikzpicture
        \let\endOrigTikzPicture\endtikzpicture
        \renewenvironment{tikzpicture}[1][]{%
          \begin{lrbox}{\SlideTikzBox}%
            \OrigTikzPicture[##1]%
          }{%
          \endOrigTikzPicture%
          \end{lrbox}%
          \resizebox{\SlideTikzPictureWidth}{\SlideTikzPictureHeight}{%
            \usebox{\SlideTikzBox}%
          }%
        }%
        \BODY
      }%
      \end{tcolorbox}%
    };
  \end{tikzpicture}%
}

% ── ពាក្យបញ្ជាជំនួយ (Utility Commands) ───────────────────────────────
\def\circF#1{}
\def\circT#1{(#1)}
\newcommand{\hoac}[1]{\left[\begin{aligned}#1\end{aligned}\right.}
\newcommand{\heva}[1]{\left\{\begin{aligned}#1\end{aligned}\right.}
\def\cautn{}
\def\cauds{}
\def\caukq{}
\def\cautl{}
\newenvironment{name}[2]{}{}
\renewcommand{\loigiai}[2][]{}        % លាក់ដំណោះស្រាយនៅលើស្លាយ
\newenvironment{indapan}[2]{}{}
\renewcommand{\shortans}[2][]{\begin{flushright}\scalebox{2.0}{$\square\square\square\square$}\end{flushright}}	
\renewcommand{\baselinestretch}{1.4}

\begin{document}

\setcounter{ex}{0}
% __SLIDE_BODY__

\end{document}
